% ================================ A.2
\subsection{\rev{zg}{}{Traffic Speed Prediction and Error Estimation}}\label{A2}

To generate realistic stochastic prediction errors, we employ the Diffusion Convolutional Recurrent Neural Network (DCRNN) \citep{li2017diffusion} to model spatiotemporal traffic dependencies. Two benchmark datasets are utilized: \textbf{METR-LA} (4 months of highway data from 207 sensors) as a proxy for Mandl’s network, and \textbf{Q-Traffic} (61 days of Beijing 6th Ring Road segments) for the real-world case study. We split both datasets, chronologically,  as recommended by \cite{li2017diffusion} and \cite{lu2022graph}, respectively: We train each prediction model on the first part of the chronologically ordered
data and measure prediction errors on the last part, i.e., the test set.
Table~\ref{tab:prediction_error} reports the one-hour-ahead prediction errors
used in our experiments.

% ----------------- table_predERROR
\begin{table}[htbp!]
\centering
\caption{One-hour-ahead prediction error on benchmark datasets.}
\label{tab:prediction_error}
\begin{tabular}{llcc}
\hline
Dataset & Model & MAE & RMSE \\
\hline
\multirow{2}{*}{METR-LA}
& DCRNN & 3.590 & 7.607 \\
& MTGNN & 3.490 & 7.230 \\
\hline
\multirow{2}{*}{Q-Traffic}
& DCRNN & 3.932 & 6.030 \\
& MTGNN & 2.973 & 4.742 \\
\hline
\end{tabular}
\end{table}
% ----------------- table_predERRORend

The prediction models output traffic speeds. Since the downstream public-transport redesign model operates on travel times, the extracted prediction residuals must be expressed in travel-time space. If a link length $d$ is available, the realized and predicted travel times are $T=d/v$ and $\hat T=d/\hat v$, respectively. 
%
The corresponding log-relative travel-time residual is therefore $\eta=\log(\hat T)-\log(T)$. For a fixed link length, this is equivalently written as $\eta=\log(v)-\log(\hat v)$, since the length term cancels out algebraically.
% ~\footnote{A short derivation is provided in App.~\ref{app:speed_to_time_log_residual}.} 
%
This equivalence is particularly useful for \textbf{METR-LA}, where reliable segment lengths are not directly available for all sensor pairs. In this case, we compute the travel-time log-relative residual directly from the observed and predicted speeds. For \textbf{Q-Traffic}, where segment lengths are explicitly provided, the same residual can be obtained either after converting speeds into travel times or directly through the equivalent speed-based expression. The resulting residual pool is then used to generate noisy travel-time forecasts as described in App.~\ref{A3}.



% -----------------------------
% -----------------------------
% -----------------------------
% ================================ A.3
\subsection{\rev{zg}{}{Empirical Log-Relative Residual Sampling for Travel-Time Prediction}}\label{A3}

We pool all empirical prediction residuals on the test sets, across links and time steps, to construct a representative empirical error pool. Since the downstream experiments transfer prediction errors across links and networks with different travel-time scales, we sample log-relative residuals rather than absolute residuals. Specifically, for each realized travel time $T$ and its prediction $\hat T$, we define $\eta=\log(\hat T)-\log(T)=\log\left(\hat T/T\right)$.
%
This log-ratio formulation is appropriate for heteroscedastic prediction errors whose magnitude scales with the level of the predicted variable~\citep{Tofallis2015}; in our setting, this is consistent with traffic scenarios in which the coefficient of variation changes the scale of realized travel times. 
% 适合误差随变量尺度增长的 heteroscedastic 情况
%
Finally, independent samples $\eta_{ij}^{\omega}$ are drawn with replacement from this empirical pool for each link $(i,j)$ and scenario $\omega$. Given the realized travel time $t_{ij}^{\omega}$ generated under the prescribed traffic-variability level, the corresponding additive travel-time error is defined as $\epsilon_{ij}^{\omega}=t_{ij}^{\omega}\left(\exp\!\left(\alpha\eta_{ij}^{\omega}\right)-1\right)$. The predicted link travel times~$\hat t_{ij}^{\omega}$, used to make network-design decisions, are then computed by adding $\epsilon_{ij}^{\omega}$ to the actual travel time~$t_{ij}^{\omega}$, as $\hat t_{ij}^{\omega}=t_{ij}^{\omega}+\epsilon_{ij}^{\omega}$. We do this in both Mandl and Beijing scenarios.

Moreover, because errors are sampled independently for each link and scenario, this first-order model abstracts away the spatial and temporal correlation of real forecast errors, such as corridor-level or incident-level common shocks. The reported results should therefore be read as an evaluation under sampled independent empirical residual scenarios, and correlation-aware forecast-error modeling remains an important extension. \rev{zg}{}{Because the simulated forecasts are constructed from realized times, this controlled error experiment does not itself validate a live, ex-ante forecasting system.}

\begin{equation}
\hat t_{ij}^{\omega} = t_{ij}^{\omega}
+
\epsilon_{ij}^{\omega}, \
%
\epsilon_{ij}^{\omega} =
t_{ij}^{\omega}
\left(
\exp\!\left(\alpha\eta_{ij}^{\omega}\right)-1
\right)
\label{eq:predTRAVEL_appendix}
\end{equation}


% ----------------------

\paragraph{Diagnostic check on traffic variability.}
For the Mandl experiments, we further report the realized-to-nominal travel-time ratio
$R_{ij}^{\omega}=t_{ij}^{\omega}/t_{ij}^{0}$, where $t_{ij}^{0}$ is the nominal link travel time and $t_{ij}^{\omega}$ is the realized travel time under scenario $\omega$. Table~\ref{tab:travel_time_ratio} summarizes the distribution of $R$ under different traffic-variability levels. The results show that increasing \textit{cv} mainly amplifies the dispersion and upper tail of realized travel times. Small-\textit{cv} scenarios remain close to nominal conditions, whereas large-\textit{cv} scenarios generate severe congested realizations, as reflected by the increasing standard deviation and 95th percentile of $R$. In these diagnostics, the prediction-error scale is fixed at $\alpha=1.0$, corresponding to the empirical error magnitude extracted from the trained prediction model.

% -------------------- table_dig
\begin{table}[htbp!]
\centering
\caption{Distribution of realized-to-nominal travel-time ratios in the Mandl experiments.}
\label{tab:travel_time_ratio}
\begin{tabular}{lccc}
\hline
\textit{cv} & Mean($R$) & Std($R$) & P95($R$) \\
\hline
0.00 & 1.000 & 0.000 & 1.000 \\
0.05 & 0.988 & 0.047 & 1.076 \\
0.25 & 0.961 & 0.228 & 1.436 \\
0.50 & 0.973 & 0.456 & 2.001 \\
1.00 & 1.090 & 0.949 & 3.396 \\
\hline
\end{tabular}
\end{table}
% -------------------- table_digEND


% % ================================ A.4
% \subsection{Equivalence Between Speed-Based and Travel-Time Log Residuals}
% \label{app:speed_to_time_log_residual}

% For a link with fixed length $d>0$, let $v$ and $\hat v$ denote the realized
% and predicted speeds, respectively. The corresponding realized and predicted
% travel times are
% $T=d/v$ and $\hat T=d/\hat v$. The log-relative travel-time residual is then
% \begin{equation}
% \begin{aligned}
% \eta
% &=
% \log(\hat T)-\log(T)
% &=
% \log\left(\frac{d}{\hat v}\right)
% -
% \log\left(\frac{d}{v}\right) 
% &=
% \left(\log d-\log \hat v\right)
% -
% \left(\log d-\log v\right) 
% &=
% \log v-\log \hat v.
% \end{aligned}
% \label{eq:speed_time_log_residual_equivalence}
% \end{equation}
